This section details the computation of the weighting factor $\norm{ \frac{\delta \tv}{\delta x} \times \frac{\delta \tv}{\delta \phi} }$, required for the calculation of the radio map due to diffraction as described in Section~\ref{sec:radio-map-solver-diffraction}.
Recall that $\tv(x, \phi)$ denotes the reparametrization of a point on the measurement cell $M_i$ reached by a diffracted ray originating from the diffraction point $\vv$ on the edge $\Ec$, with $\phi$ representing the Keller cone azimuth.

We begin by rewriting~\eqref{eq:radio-map-solver-ci-D-reparametrization} as
\begin{equation}
    \label{eq:diffraction-radio-map-weighting-factor-reparametrization}
    \tv(x,\phi) = \vv(x) + \gamma \widehat{\kv}_s(x, \phi),
\end{equation}
where $\vv(x) = \ov + x \widehat{\ev}$ is the diffraction point along the edge $\Ec$, and $\widehat{\kv}_s$ denotes the direction of the diffracted ray, defined by the Keller cone azimuth $\phi$ and the opening angle $\beta_0'$:
\begin{equation}
    \label{eq:diffraction-radio-map-weighting-factor-kv-s}
    \widehat{\kv}_s(x, \phi) = \sin\beta_0' \cos\phi\, \widehat{\tv}_0 + \sin\beta_0' \sin\phi\, \widehat{\nv}_0 + \cos\beta_0'\, \widehat{\ev}.
\end{equation}
The dependencies of $\vv$ and $\widehat{\kv}_s$ on $x$ and $\phi$ are shown explicitly for clarity.
Let $\widehat{\nv}$ denote the normal to the measurement cell $M_i$, and $K$ a point on the plane containing the measurement cell.
Since $\tv$ lies on $M_i$, $\gamma$ can be found by requiring
\begin{equation}
    \nv\tp \LB \tv(x, \phi) - K \RB = 0
\end{equation}
which yields
\begin{equation}
    \label{eq:diffraction-radio-map-weighting-factor-gamma}
    \gamma = \frac{\nv\tp \LB K - \vv(x) \RB}{\nv\tp \widehat{\kv}_s(x, \phi)}.
\end{equation}
Additionally, for a fixed source position $\sv$, the angle of incidence $\beta_0'$ depends on the location of the diffraction point $\vv(x)$ along edge $\Ec$
\begin{equation}
    \begin{aligned}
        \cos\beta_0' &= \frac{\norm{\vv(x) - \sv_{\Ec}}}{\norm{\vv(x) - \sv}}, \\
        \sin\beta_0' &= \frac{\norm{\sv - \sv_{\Ec}}}{\norm{\vv(x) - \sv}}
    \end{aligned}
    \label{eq:diffraction-radio-map-weighting-factor-beta-0}
\end{equation}
where $\sv_{\Ec}$ denotes the projection of the source point $\sv$ onto the edge $\Ec$.

Equations~\eqref{eq:diffraction-radio-map-weighting-factor-reparametrization}, \eqref{eq:diffraction-radio-map-weighting-factor-kv-s}, \eqref{eq:diffraction-radio-map-weighting-factor-gamma}, and \eqref{eq:diffraction-radio-map-weighting-factor-beta-0} together provide the reparametrization of $\tv$ in terms of $x$ and $\phi$, for a given source position $\sv$.
In \SRT{}, the derivatives of $\tv$ with respect to $x$ and $\phi$ are computed via automatic differentiation.
Specifically, Dr.Jit~\cite{Jakob2022DrJit} is employed to obtain $\frac{\delta \tv}{\delta x} \in \RR^3$ and $\frac{\delta \tv}{\delta \phi} \in \RR^3$ by differentiating through these equations. The weighting factor is then evaluated as the norm of the cross product of these two derivatives.
One could also have implemented the derivatives by hand, but using automatic differentiation simplifies the implementation without incurring significant overhead.